\textbf{Proof.}
Put \(L=\log(en)\).  We prove two bands and choose the numerical constants at the end.

\emph{Paired-sign band.}
Let \(k=2\lfloor d/2\rfloor\), let \(\mathbf Q\) be uniform on its first \(k\) coordinates and zero elsewhere, and, for \(\sigma\in\{-1,1\}^{k/2}\), define
\[
 P_{\sigma,2j-1}=\frac{1+\alpha\sigma_j}{k},
 \qquad
 P_{\sigma,2j}=\frac{1-\alpha\sigma_j}{k},
 \qquad
 \alpha^2=\min\!\left\{\frac1{16},\frac{\sqrt{k}}{16n}\right\}.
\]
Then \(0<\alpha\leq1/4\), so every \(\mathbf P_\sigma\) is a probability vector.  Direct summation gives
\[
 \lVert\mathbf P_\sigma-\mathbf Q\rVert_1=\alpha.
\]
Let \(M\) be the uniform mixture of the laws \(\mathbf P_\sigma^n\).  If \(\sigma\) and \(\tau\) are independent uniform sign vectors, the likelihood-ratio second-moment identity gives
\[
 \begin{aligned}
 1+\chi^2(M,\mathbf Q^n)
 &=E_{\sigma,\tau}\left[
   1+\frac{2\alpha^2}{k}
          \sum_{j=1}^{k/2}\sigma_j\tau_j
   \right]^n\\
 &\leq E_{\sigma,\tau}\exp\!\left{
   \frac{2n\alpha^2}{k}
          \sum_{j=1}^{k/2}\sigma_j\tau_j
   \right}\\
 &=\left\{\cosh\!\left(\frac{2n\alpha^2}{k}\right)\right\}^{k/2}
 \leq\exp\!\left(\frac{n^2\alpha^4}{k}\right)
 \leq e^{1/256}.
 \end{aligned}
\]
The last inequality holds with equality in the exponent when \(\alpha^2=\sqrt{k}/(16n)\); when \(\alpha^2=1/16\), the condition selecting the cap gives \(k\geq n^2\).  The other inequalities use \(1+t\leq e^t\) and \(\cosh t\leq e^{t^2/2}\).  Consequently,
\[
 \lVert M-\mathbf Q^n\rVert_{\mathrm{TV}}
 \leq\frac12\sqrt{e^{1/256}-1}<\frac12.
\]
For any estimator \(\widehat L\), test the null target zero against the mixture target \(\alpha\) by \(\phi=\mathbf 1\{\widehat L\geq\alpha/2\}\).  The maximum risk is at least one half of the sum of its risk under the null and its mixture-average risk under the alternatives.  Hence
\[
 \begin{aligned}
 \sup_{(\mathbf P,\mathbf Q)\in\Delta_d^2}
 E\!\left[(\widehat L-\lVert\mathbf P-\mathbf Q\rVert_1)^2\right]
 &\geq\frac{\alpha^2}{8}
   \left\{\mathbf Q^n(\phi=1)+M(\phi=0)\right\}\\
 &\geq\frac{\alpha^2}{8}
   \left\{1-\lVert M-\mathbf Q^n\rVert_{\mathrm{TV}}\right\}
 \geq\frac{\alpha^2}{16}.
 \end{aligned}
\]
The independent sample from the same fixed \(\mathbf Q\) under every hypothesis contains no information and leaves this testing bound unchanged.  Taking the infimum over \(\widehat L\) gives
\[
 \mathfrak R^{L_1}_{n,d}
 \geq\frac1{256}\min\!\left\{1,\frac{\sqrt{k}}n\right\}.
\]
If \(2\leq d\leq DL^2\), then \(k\geq d/2\).  On \(\sqrt{k}\leq n\), the last display and \(\sqrt d/L\leq\sqrt D\) give a constant multiple of \(d/(nL)\).  On \(\sqrt{k}>n\), its right-hand side is \(1/256\), while \(d/(nL)\leq DL/n\leq D\).  Thus, after decreasing a constant depending only on the fixed numerical \(D\),
\[
 \mathfrak R^{L_1}_{n,d}\geq c_D\frac{d}{nL},
 \qquad 2\leq d\leq DL^2,
\]
without an additional range restriction.

\emph{Uniform known-\(\mathbf Q\) band.}
Now take \(Q_x=1/d\) for every \(x\).  Fix the logarithmic constant in \(\mathrm{lem:jiao-han-weissman-known-q-lower}\) to be \(1/2\), and denote its corresponding dominance constant by \(C'>0\).  For all sufficiently large \(n\), the restriction \(d\leq\rho nL\), with a sufficiently small numerical \(\rho\), implies
\[
 \log n\geq\frac12\log d,
 \qquad
 T_{n,d}(\mathbf Q)=\sqrt{\frac{d}{n\log n}}.
\]
If \(d\geq DL^2\), then
\[
 \frac{T_{n,d}(\mathbf Q)}{\sqrt{\log n/n}}
 =\frac{\sqrt d}{\log n}
 \geq\sqrt D,
 \qquad
 \frac{T_{n,d}(\mathbf Q)}{\sqrt d\log n/n}
 =\frac{\sqrt n}{(\log n)^{3/2}}.
\]
Choose \(D\) so that the first ratio is at least \(2C'\), and then choose \(n_0\) so that the second ratio is at least \(2C'\) for \(n\geq n_0\).  Enlarging \(D\) once more makes the inequalities \(d\geq DL^2\) and \(d\leq\rho nL\) incompatible for the finitely many \(2\leq n<n_0\).  Thus the written dominance condition of the cited lemma holds throughout the nonempty upper band.  Its fixed-multinomial conclusion gives
\[
 \inf_{\widehat L_{\mathbf Q}}
 \sup_{\mathbf P\in\Delta_d}
 E_{\mathbf P}\!\left[
   (\widehat L_{\mathbf Q}-\lVert\mathbf P-\mathbf Q\rVert_1)^2
 \right]
 \geq c_0\frac{d}{n\log n}
 \geq c_1\frac{d}{nL}.
\]
Here the estimator on the left knows \(\mathbf Q\).  Given any estimator in the unknown--unknown two-sample experiment, an estimator in the known-\(\mathbf Q\) experiment can simulate the missing independent \(\mathbf Q\)-sample and apply it.  Conditional expectation given the \(\mathbf P\)-sample derandomizes this procedure without increasing squared risk.  Therefore the unknown--unknown minimax risk is no smaller than the displayed known-\(\mathbf Q\) risk.

The paired-sign argument covers \(d\leq DL^2\), while the cited argument covers \(d\geq DL^2\).  Taking the smaller of their positive constants and the chosen sufficiently small \(\rho\) proves the claim.  All support constraints were checked above; no boundary coordinate with zero \(Q_x\) is used in the cited uniform band.